\documentclass{article}
\usepackage{amsmath,amssymb}
\usepackage{xurl}
\usepackage[hidelinks]{hyperref}
% Shared commands for notes in docs/paper-gaps/.

\DeclareUnicodeCharacter{2080}{\ifmmode{}_{0}\else\textsubscript{0}\fi}
\DeclareUnicodeCharacter{2081}{\ifmmode{}_{1}\else\textsubscript{1}\fi}
\DeclareUnicodeCharacter{2082}{\ifmmode{}_{2}\else\textsubscript{2}\fi}
\DeclareUnicodeCharacter{2099}{\ifmmode{}_{n}\else\textsubscript{n}\fi}

\newcommand{\Lean}{\textsc{Lean}}
\ExplSyntaxOn
\NewDocumentCommand{\leanid}{m}{
  \ifmmode
    \textsf{\detokenize{#1}}
  \else
    \group_begin:
    \urlstyle{sf}
    \tl_set:Nn \l_tmpa_tl {#1}
    \tl_replace_all:Nnn \l_tmpa_tl {\_} {_}
    \exp_args:NV \path \l_tmpa_tl
    \group_end:
  \fi
}
\ExplSyntaxOff
\providecommand{\tr}{\operatorname{tr}}
\newcommand{\tnleanIssueBase}{https://github.com/LionSR/TNLean/issues/}
\newcommand{\tnleanPrBase}{https://github.com/LionSR/TNLean/pull/}
\newcommand{\tnleanDocBase}{https://sirui-lu.com/TNLean/docs/}
\newcommand{\ghissue}[1]{\href{\tnleanIssueBase#1}{GitHub issue \##1}}
\newcommand{\ghpr}[1]{\href{\tnleanPrBase#1}{PR \##1}}
\newcommand{\doclink}[1]{\href{\tnleanDocBase#1.html}{\path{#1}}}

% Dirac notation and the identity operator, as in blueprint/src/macros/common.tex.
% \providecommand keeps note-local definitions authoritative.
\usepackage{dsfont}
\providecommand{\ket}[1]{|#1\rangle}
\providecommand{\bra}[1]{\langle #1|}
\providecommand{\braket}[2]{\langle #1 | #2 \rangle}
\providecommand{\ketbra}[2]{|#1\rangle\!\langle #2|}
\providecommand{\Id}{\mathds{1}}
\providecommand{\one}{\mathds{1}}
\providecommand{\C}{\mathbb{C}}
\providecommand{\MN}[1]{M_{#1}(\C)}
\providecommand{\MD}{M_D(\C)}

% Verdict marker: \gapnote{<kind>}{<status>}. Kind is the policy
% classification (clarification, local-correction, scope-restriction,
% unfaithful, false-source, open-gap); status is open, wip (elimination
% underway), resolved, or historical. Machine-read by the site index;
% prints a verdict line.
\newcommand{\gapnote}[2]{\par\noindent\textbf{Verdict:} #1 (#2).\par}

\title{Circuits Between States in the Same Phase: the Normal Case}
\author{}
\date{2026-10-01}
\begin{document}
\maketitle
\gapnote{scope-restriction}{resolved}

\section{What the source states}
The paragraph ``Discussion and outlook'' of~\cite{Malz2024Preparation} states:
``Our results also imply that MPS in the same phase can be transformed into
each other using a log-depth circuit''. No proof is given. The results it
refers to are the preparation of normal states in depth
$O(\log(N/\varepsilon))$, eq.~(1), and the preparation of all
translation-invariant states, ``including long-range non-normal ones'', with
measurements. The main text recalls that ``normal MPS all lie in the
topologically trivial phase''.

\section{What is formalized}
Let $A$ and $B$ be normal tensors with the same physical dimension, and
write $\hat\phi_N(A)=\ket{\phi_N(A)}/\lVert\phi_N(A)\rVert$ for the normalized periodic
states and
$\varepsilon(\psi,\phi)=1-\lvert\braket{\psi}{\phi}\rvert$ for the error of
eq.~(1). There is $c$, depending only on $A$ and $B$, such that for $N\ge2$,
$0<\varepsilon\le1$ and nonvanishing periodic states, a local circuit $U$ of
depth at most $c\log(N/\varepsilon)$ satisfies
$\varepsilon(U\hat\phi_N(A),\hat\phi_N(B))\le\varepsilon$; if both bond dimensions are
positive, the same holds for all $N\ge N_0$ without the nonvanishing
condition.\footnote{Formal declarations:
    \leanid{MPSPreparation.exists_isLocalCircuitOfDepth_le_log_of_mpvState_ne_zero}
    and \leanid{MPSPreparation.exists_isLocalCircuitOfDepth_le_log} in
    \doclink{TNLean/MPS/Preparation/CircuitEquivalence}.}
The circuit is $U_BU_A^\dagger$, with $U_A$ and $U_B$ the preparations of
eq.~(1) from the all-zero product state, each with error $\varepsilon/4$. The
error is not a metric; the proof uses
$\varepsilon(x,z)\le2\bigl(\varepsilon(x,y)+\varepsilon(y,z)\bigr)$ for unit
vectors, since $2\varepsilon(x,y)$ is the squared distance from $x$ to the
nearest unit multiple of $y$.

\section{The restriction}
The theorem above covers the trivial phase of the classification without
symmetry: both tensors are normal. Phases defined with a symmetry would need
circuits that commute with the symmetry, which the source does not claim, and
are not treated.

\section{The non-normal case}
For a non-normal tensor the parent Hamiltonian has a degenerate ground space,
spanned by the periodic states of the blocks of the canonical form, and two
such tensors lie in the same phase without symmetry exactly when they have the
same number of blocks. The source's own construction of long-range states,
eq.~(25), prepares the block states together with length-dependent
coefficients. The object that a circuit between states of the same phase
transports is therefore the ground space, and that statement is formalized:
for two matched families of $b$ normal blocks, one nearest-neighbour unitary
of depth $O(\log(N/\varepsilon))$ maps the fixed encoder of one ground space
to that of the other within $\varepsilon$ in operator norm, uniformly over
logical inputs and finite references; see
\cite{gap:mswc24_coherent_sector_encoder_scope}.\footnote{Formal declaration:
    \leanid{MPSPreparation.exists_log_depth_coherent_sectorEncoder_conversion}.}

The stronger reading, in which one chosen normalized periodic vector is mapped
to another chosen vector, is false: two tensors with the same parent
Hamiltonian but sector probabilities $(4/5,1/5)$ and $(1/2,1/2)$ cannot be
converted by a nearest-neighbour circuit of depth $T$ with infidelity below
$9/5000$ once $N>4T+4$.\footnote{Formal declarations:
    \leanid{MPSTensor.parentHamiltonian_repeatedBlockTensor_eq_ghzTensor} and
    \leanid{MPSTensor.repeatedBlockTensor_infidelity_lower_bound}.}
Equality of phases therefore does not transport sector probabilities, and the
ground-space statement is the faithful reading of the source's remark.
Compatible coherent seeds give a further sufficient condition for converting
chosen vectors; see \cite{gap:mswc24_same_phase_compatible_seed}.

\bibliographystyle{alpha}
\bibliography{references}
\end{document}
